\chapter{Die Notfall-KI: Ein Sprachassistent ohne Internet}

NOMAD kann einen Sprachassistenten (KI) betreiben, der \textbf{auf Ihrem Rechner} läuft: Sie stellen Fragen in ganzen Sätzen, die Antwort wird lokal berechnet. Es geht nichts ins Internet.

\kasten{warnrot}{Was die KI nicht ist}{%
Die KI kann sich irren und Dinge erfinden. Bei Medizin, Sicherheit, Recht und Chemikalien \textbf{prüfen Sie jede Antwort an den Quellen} in der Wissensbibliothek (Kapitel~5) oder in gedruckten Unterlagen. Sie ist ein Helfer zum Suchen und Erklären, kein Arzt.}

\section{Was dafür nötig ist}
\begin{itemize}[leftmargin=1.5em,itemsep=3pt]
\item \textbf{Internet beim Einrichten:} Das Sprachmodell (je nach Größe 1 bis mehrere GB) wird einmal heruntergeladen.
\item \textbf{Arbeitsspeicher:} Je größer das Modell, desto besser die Antworten, aber desto mehr Speicher und Zeit. Richtwerte*:
\end{itemize}

\begin{center}\small
\renewcommand{\arraystretch}{1.3}
\begin{tabularx}{\linewidth}{@{}>{\RaggedRight\arraybackslash}p{4.2cm}>{\RaggedRight\arraybackslash}p{3.5cm}X@{}}
\toprule
\textbf{Rechner} & \textbf{Modellgröße*} & \textbf{Zu erwarten*} \\
\midrule
Wenig Speicher (4--8~GB), keine Grafikkarte & ca. 1--3 Milliarden Parameter & einfache Fragen; langsam bis ordentlich \\
\midrule
Mittel (16~GB) & ca. 7--9 Milliarden & brauchbare deutsche Antworten, bei Rechnern ohne Grafikkarte einige Sekunden bis Minuten pro Antwort \\
\midrule
Grafikkarte (NVIDIA, 8~GB+ Speicher) & ca. 8--14 Milliarden & flott, bessere Qualität \\
\bottomrule
\end{tabularx}
\end{center}
\small * Diese Richtwerte stammen aus allgemeiner Erfahrung und sind von uns für dieses Heft noch \textbf{nicht} auf Ihrem Rechner gemessen. NOMAD empfiehlt selbst, mit kleinen Modellen zu beginnen und größere auszuprobieren.\normalsize

\section{Einrichten}
\begin{enumerate}[leftmargin=1.5em,itemsep=3pt]
\item Internet anschließen.
\item Startseite $\rightarrow$ \menue{Schnellstart} oder \menue{Einstellungen} $\rightarrow$ \menue{Apps verwalten}: den \textbf{KI-Assistenten} installieren.
\item In den KI-Einstellungen ein \textbf{kleines} Modell wählen und laden. Erst probieren, dann ein größeres.
\item Internet trennen, Frage stellen (Kapitel~6, Checkliste).
\end{enumerate}

\hinweis{Hat der Rechner eine NVIDIA-Grafikkarte, nutzt NOMAD sie, wenn der Installer sie erkennt (steht am Ende der Installation). Ohne Grafikkarte läuft die KI auf dem Prozessor und ist langsamer.}

\section{Eigene Dokumente für die KI}
In der Wissensdatenbank können Sie Dateien (PDF, Text) hochladen, aus denen die KI Antworten ableitet. Sinnvoll sind Ihre eigenen Notfallpläne, Telefonlisten, das Handbuch als PDF oder Medikamentenpläne der Familie.

\kasten{hinweisblau}{Noch nicht geprüft}{%
Ob die KI auch die per Kiwix geladenen Bücher (Wikipedia, Nachschlagewerke) \emph{selbst} nutzt, haben wir noch nicht abschließend geprüft. Verlassen Sie sich nicht darauf: Suchen Sie bei wichtigen Fragen selbst in der Wissensbibliothek.}

\section{Fragen stellen}
\begin{itemize}[leftmargin=1.5em,itemsep=3pt]
\item Kurz und konkret fragen: „Wie bereite ich Wasser mit Tabletten auf?“
\item Bei Zahlen und Dosierungen immer gegenprüfen.
\item Antwort unklar? Frage anders formulieren oder die Quelle in der Wissensbibliothek lesen.
\end{itemize}

\section{Ihre Notizen zur KI}
\renewcommand{\arraystretch}{1.9}
\begin{tabularx}{\linewidth}{@{}>{\RaggedRight\arraybackslash}p{4.8cm}X@{}}
\toprule
Gewähltes Modell & \\
\midrule
Größe auf der Platte & \\
\midrule
Dauer einer typischen Antwort & \\
\midrule
Hochgeladene Dokumente & \\
\bottomrule
\end{tabularx}
